\section{Discussion}
\label{sec:discussion}

\subsection{Deployment recommendations}
The results support a concise set of recommendations. In every case, use the plain question, one frame per decision, and a threshold on $p_\text{yes}$ chosen on labelled footage from the site rather than the model's own yes/no answer.

\textit{Device with an integrated GPU or accelerator.} LFM2.5-VL-1.6B with Q4\_K\_M weights and 256\,px input offers the best balance we measured: F1 \CLfmMQfourTwofiftysixOneFone, recall \CLfmMQfourTwofiftysixOneRecall\%, hard-negative false alarms \CLfmMQfourTwofiftysixOneFarHard\%, p95 latency \CLfmMQfourTwofiftysixOnePtf\,ms and \CLfmMQfourTwofiftysixOneRss\,GiB, sufficient for \CLfmMQfourTwofiftysixOneStreams{} cameras at one decision per second each, at \PLfmMQfourTwofiftysixOneGpuJ\,J per decision. Where accuracy outweighs camera count, LFM2.5-VL-3B with Q8\_0 weights at 512\,px reaches F1 \CLfmLQeightFivetwelveOneFone{} with \CLfmLQeightFivetwelveOneFarHard\% hard-negative false alarms, at \CLfmLQeightFivetwelveOnePtf\,ms and \CLfmLQeightFivetwelveOneRss\,GiB.

\textit{CPU only.} Only LFM2.5-VL-450M sustains one decision per second for one camera at useful accuracy: with Q8\_0 weights at 512\,px it reaches F1 \CLfmSQeightFivetwelveOneFone{} (recall \CLfmSQeightFivetwelveOneRecall\%, hard-negative false alarms \CLfmSQeightFivetwelveOneFarHard\%) at \CLfmSQeightFivetwelveOneCpuPtf\,ms, \CLfmSQeightFivetwelveOneRss\,GiB and \PLfmSQeightFivetwelveOneCpuJ\,J per decision. If a slower check rate is tolerable, LFM2.5-VL-1.6B at 256\,px gives F1 \CLfmMQfourTwofiftysixOneFone{} at one decision every \CLfmMQfourTwofiftysixOneCpuPtf\,ms per camera. Our CPU is a fast laptop core; slower embedded CPUs will need the 450\,M model or a lower decision rate. A GPU or accelerator, where available, should be used: the same configuration consumed \PCpuGpuRatioMin--\PCpuGpuRatioMax{} times more energy per decision on the CPU.

\textit{What to avoid.} Long, definition-style prompts; 256\,px input with any model but LFM2.5-VL-1.6B; multi-frame inputs, which multiply latency without improving accuracy; and 16-bit weights, which occupy about twice the memory of Q4\_K\_M with no consistent accuracy gain.

\subsection{Prompting small models}
A prompt that defines the task meticulously and enumerates the cases to exclude reads better to a human, yet in our grid it drove the 450\,M model to answer ``yes'' to nearly every frame and raised false alarms for every LFM model (Table~\ref{tab:prompts}). VLMs are known to answer ``yes'' too readily \cite{li2023pope,chan2026yesbias}, and in thermal fall screening a single cautionary sentence shifted sensitivity substantially \cite{junca2026thermal}. Our results suggest that naming the excluded postures amplifies the yes-bias of small models instead of correcting it. A one-sentence criterion (\promptBtwo) looked best for Qwen3.5-2B in the 30-frame pilot, but on the full benchmark it rendered that model conservative, and it missed about half of the falls. Across all five models the plain question was best or tied. The pilot is a salutary reminder of how misleading 30 samples can be.

\subsection{Thresholds instead of answers}
The default answer is the weakest element of these models. Their scores rank falls above look-alike postures well (AUROC \RLfmSAuroc--\RLfmLAuroc), but the cut-off implied by the yes/no answer lies elsewhere for each model: LFM2.5-VL-450M answered ``yes'' for \RLfmSDefFarHard\% of hard negatives, Qwen3.5-4B for \RQwenLDefFarHard\%. Choosing the threshold on development data reduced false alarms for every model, most for the smallest. Token probabilities are not necessarily calibrated \cite{guo2017calibration}, so the threshold must be learned from labelled data, and it does not transfer perfectly: two models reached 90--91\% recall on the test set with thresholds set for 95\% on development. Reading token probabilities is free in llama.cpp and most runtimes, so deployments should log $p_\text{yes}$ and set, and periodically re-check, a site-specific threshold.

\subsection{Is a VLM needed?}
For this task, yes. The pose-and-rule baseline reached F1 \CPoseNPoseNatOneFone--\CPoseNPoseNatFourFone, against at least \RLfmSCalFone{} for every VLM, because the pose detector frequently fails to find a person lying on the floor at all. Its speed (\RPoseNMs\,ms per frame) is attractive, but an inexpensive first stage that overlooks about half of the fallen people cannot gate a slower model. A pose model trained on fallen people, or a detector-guided prompt as in \cite{adil2026svlm}, might alter this verdict; we leave both to future work.

\subsection{Measuring efficiency}
Two details of our procedure changed results by large factors. The first is memory accounting on unified-memory hardware: depending on the metric and on two runtime flags, the same 3\,B model ``used'' between \AblLargeGpuDefaultFpGiB\ and \AblLargeCpuDefaultRssGiB\,GiB (Table~\ref{tab:memory}). Papers that report one peak-memory figure without naming the metric and runtime settings resist comparison. Weight repacking alone accounted for about 2.5\,GiB of the 3\,B model's CPU memory, which matters on an 8\,GB board. The second is device state: in the pilot, a battery run drifted upward in latency by 15--23\% over about 30 minutes, and idle sleep stalled one request for 17 minutes; during the main grid, 41 minutes on battery slowed the affected configurations until they were re-measured. Neither problem is peculiar to our setup, and both are easy to overlook without per-request logs and a power log. Answers, by contrast, were fully reproducible: re-running \RRetimeConfigs{} configurations reproduced \RRetimeSame{} of \RRetimeTotal{} decisions.

\subsection{Limitations}
\label{sec:limitations}
\textit{Domain gap.} Both datasets are indoor, staged recordings in homes, offices and a laboratory, with clean floors, adequate lighting and few occluders. Industrial floors add machinery, pallets, high-visibility clothing, hard hats and partial occlusion. We did not test on industrial footage, so our figures are best read as an upper bound on factory performance until confirmed there.

\textit{Two datasets, few actors, staged falls.} MCF has a single actor and UR Fall a few; falls are performed, and in MCF mostly onto a mattress, which may itself serve as a cue. Real collapses look different. UR Fall's overhead camera recorded only falls, so overhead false-alarm rates cannot be measured. Some hard-negative types are small (16 kneeling and 22 crouching test samples), as the confidence intervals reveal.

\textit{One device.} All efficiency figures come from one laptop-class system-on-chip, used as a CPU-only and an integrated-GPU operating point. Embedded boards will be slower and have less memory; relative effects of size, precision, resolution and frame count, and the device-independent costs we report, should transfer better than absolute latencies. Power is the package estimate from \texttt{powermetrics}, which excludes memory and other components, so wall-power figures would be higher.

\textit{Single decisions.} We evaluate single frames or short clips. A deployed system would aggregate decisions over time, for example by requiring several consecutive positives, which curbs false alarms at the cost of delay. We do not model this.

\textit{Zero-shot only.} We fine-tune nothing. Fine-tuning would probably raise accuracy, especially for the smallest models, but demands target-domain training data.

\textit{Labels and thresholds.} Labels were reviewed by an AI assistant, then checked by one human, blind on a sample and with labels visible on all frames; \HumanCheckText{} Independent human labelling of every frame would strengthen the benchmark, and subtype labels (crouching versus kneeling, say) are the most uncertain. The calibrated operating point uses development data from the same datasets as the test set, so it is optimistic about how a threshold transfers to a new site.

\textit{Model snapshot and licences.} Small VLMs evolve quickly; our conclusions apply to the versions listed, and the benchmark and scripts are released so newer models can be added. UR Fall restricts use to non-commercial research and MCF states no licence, constraining what we can redistribute and possibly the commercial reuse of derived results.

\subsection{Ethical considerations}
Camera-based monitoring of workers raises privacy concerns. A \pd{} detector is meant to prompt a check by a person, not to monitor productivity or discipline workers, and local processing keeps video on site. We used only public datasets collected by others and recorded no one. In figures we blur faces and show only frames from a dataset whose licence permits it. Any deployment should comply with local law on workplace monitoring and involve workers or their representatives.
